%===============================================================================
% 中英文摘要
%===============================================================================

% ---------- 中文摘要 ----------
\newpage
\thispagestyle{plain}
\begin{center}
    \heiti\fontsize{18pt}{27pt}\selectfont 摘\hskip1em 要
\end{center}
\vspace{2.8mm}

\fontsize{12pt}{18pt}\selectfont
原子核是由质子和中子构成的致密强关联有限量子多体系统。在核素图上，当质子或中子离开双幻数闭壳时，原子核展现出从球形振动向空间形变转动的丰富几何相变。其中，非轴对称的\term{三轴形变}（Triaxial Deformation，即椭球三个内禀主轴长度各不相同）打破了转动对称轴守恒，不仅诱发了非谐振摆动运动（Wobbling Mode）、原子核手征对称性自发破缺（Nuclear Chirality）以及三轴超形变（TSD）等极具吸引力的奇特转动模式，而且与空间反射非对称的\term{八极关联}（Octupole Correlations）、同核异能态（Isomers）以及微观单极张量力驱动的壳层演化紧密交织。深入研究三轴形变核心与集体激发动力学，是检验微观核多体理论与强相互作用极端对称性的终极试金石。

然而，传统实验研究中往往存在诸多认知误区与判据混淆：例如将能量曲线的弯曲简单等同于横向摆动、在多极混合比反演中忽视代数双实根退化导致误判跃迁多极性、以及在四极形变评价中误将受限自旋子通道截断当成全空间电荷几何分布。针对上述关键问题，本学位论文以“易懂深入、透彻直观”的学术准则，建立面向原子核三轴形变与集体激发的系统性理论与实验分析体系：

首先，论文确立了\term{“以三轴形变核心动力学为基石，以八极形变与衰变/同核异能态为前沿拓展两翼”}的研究架构，跨越 $A \approx 60, 80, 100, 130, 160, 190$ 等六大三轴形变典型核区，系统梳理了高 $j$ 闯入轨道与费米面位置对三轴集体运动形态的微观决定作用，澄清了横向摆动（TW）与大振幅准粒子倾斜进动（TiP）的七维仲裁判据。

其次，深入推导了从微观壳模型、准粒子推转模型（CSM）、三轴粒子-转子模型（PRM）到三维倾斜轴推转模型（3D-TAC）的理论演化，阐明了转动坐标系下 Coriolis 力诱发的顺磁准粒子取向、能带回避交叉相互作用矩阵元 $|V_{12}|$ 对尖锐回弯与平滑上弯的微观控制，以及模型无关四极形状不变量 $q_2, q_3, \cos 3\gamma$ 的电荷闭合定理。

在实验技术与数据处理方法方面，系统解构了现代在束 $\gamma$ 谱学全流程：涵盖重离子熔合蒸发反应动力学、高纯锗探测器（HPGe）电荷收集与反康普顿 BGO 屏蔽机制、数字脉冲成形与弹道亏损消除、$\gamma$-$\gamma$ 符合矩阵扣除、基于 Yamazaki $F$ 系数的方向角关联（DCO/ADO）严格双根代数求解、康普顿散射线偏振度测定以及利用反冲距离多普勒频移法（RDDS）和多普勒频移衰减法（DSAM）提取皮秒级激发态寿命与绝对跃迁几率 $B(E2)/B(M1)$。

最后，论文总结了三轴集体激发、宇称-手征联立破缺四重带以及超重核合成动力学等遗留前沿问题，并对惠州强流重离子加速器装置（HIAF）等下一代大科学工程的前景进行了展望。全文力求扫清教科书式的公式堆砌与雷同套话，旨在为初学入门者与从事在束核谱学实验研究的科研人员提供一部概念清晰、物理图像鲜明且兼具实战指南价值的高水平参考论著。

\vspace{\baselineskip}
\noindent{\heiti\fontsize{14pt}{21pt}\selectfont 关键词：}
{\fontsize{12pt}{18pt}\selectfont 原子核结构；三轴形变；摆动运动；手征双重带；八极关联；在束 $\gamma$ 谱学；四极形状不变量}

% ---------- 英文摘要 ----------
\newpage
\thispagestyle{plain}
\begin{center}
    \bfseries\fontsize{18pt}{27pt}\selectfont Abstract
\end{center}
\vspace{2.8mm}

\fontsize{12pt}{18pt}\selectfont
The atomic nucleus is a strongly correlated, finite quantum many-body system composed of protons and neutrons. Away from shell closures, atomic nuclei exhibit rich geometric phase transitions from spherical vibration to deformed rotation. In particular, triaxial non-axial deformation ($\gamma \approx 15^\circ \sim 35^\circ$), where the three principal moments of inertia differ, breaks rotational axial symmetry. This gives rise to fascinating collective phenomena such as wobbling motion, spontaneous breaking of nuclear chirality, and triaxial superdeformation (TSD), which are further intertwined with reflection-asymmetric octupole correlations, nuclear isomerism, and monopole tensor force driven shell evolution. Investigating the core dynamics of triaxial deformation and collective modes provides a stringent testing ground for modern microscopic nuclear structure theory.

However, experimental spectroscopy has often suffered from ambiguities, such as mistaking energy staggering for transverse wobbling, neglecting algebraic dual-root degeneracy in multipole mixing ratios ($\delta$), or treating spin-channel-restricted quantities as complete electric charge distributions. Addressing these challenges with pedagogical clarity and physical rigor, this dissertation establishes a comprehensive theoretical and experimental framework for triaxial nuclear structure:

First, an integrated architecture—anchored by the triaxial collective core and expanded through octupole and weak-decay wings—is established across six prototypical mass regions ($A \approx 60, 80, 100, 130, 160, 190$). The microscopic driving mechanisms governed by high-$j$ intruder orbitals and Fermi surface positions are systematically revealed, clarifying the discriminating criteria between genuine transverse wobbling (TW) and tilted precession (TiP).

Second, the microscopic theoretical evolution from the independent particle shell model and the Cranked Shell Model (CSM) to the Particle-Rotor Model (PRM) and 3D Tilted Axis Cranking (3D-TAC) is thoroughly elucidated. We examine the Coriolis alignment of quasiparticles, bandcrossing interaction matrix elements $|V_{12}|$ controlling sharp backbending versus smooth upbending, and model-independent quadrupole shape invariants ($q_2, q_3, \cos 3\gamma$) with charge-closure constraints.

Third, the methodologies of in-beam $\gamma$-ray spectroscopy are systematically formulated: heavy-ion fusion-evaporation kinematics, high-purity germanium (HPGe) pulse collection and BGO Compton suppression, digital pulse shaping, $\gamma$-$\gamma$ coincidence analysis, rigorous dual-root DCO/ADO solving via Yamazaki $F$-coefficients, Compton linear polarization, and picosecond level lifetime extraction via RDDS and DSAM techniques.

Finally, unresolved puzzles regarding multi-phonon wobbling, parity-chiral quadruplet bands, and superheavy element synthesis are summarized, with perspectives for next-generation facilities such as the High Intensity Heavy-ion Accelerator Facility (HIAF). This work provides an accessible, rigorous, and practical treatise on low-energy nuclear spectroscopy.

\vspace{\baselineskip}
\noindent{\bfseries\fontsize{14pt}{21pt}\selectfont Key Words: }
{\fontsize{12pt}{18pt}\selectfont Nuclear Structure; Triaxial Deformation; Wobbling Motion; Chiral Doublet Bands; Octupole Correlations; In-Beam $\gamma$-Ray Spectroscopy; Quadrupole Invariants}
